% ==================================================================
% §15.5.3 REBUILT — Phase 1: Joint effective allowed band and working benchmark
%
% This block is Part 3 of the §15.5 rebuild under the parallel-creation
% strategy (PROJECT_STRUCTURE.md §3.6). Insert IMMEDIATELY AFTER
% §15.5.2 block (label: subsec:eps_retained_rebuilt) and BEFORE
% the existing old §15.5.
%
% GPT round 7 (10 corrections) + user directive 2026-04-18 applied:
%   NO legacy Hubble compensation numbers (3%, 2.85%, 2.90%, 69 km/s/Mpc)
%   NO legacy specific ε values (0.010, 0.012)
%   Midpoint ε ≈ 0.032 used ONLY as service illustrative quantity.
%   Range-based scaling table (not midpoint-based).
%   Benchmark A retained as legacy label only, with no numerics.
%   No figure label name (figure pending regeneration).
%   No reference to epsilon_diagram.html.
%
% Migration items from old §15.5 absorbed here per 02_s15_5_2_INVENTORY.md:
%   #2 (Bridge from age apparatus) → opening bridge cross-ref
%   (items #1, #11c, #13, #14b reclassified: demote-to-appendix X2 ONLY,
%    NOT migrated to §15.5.3 main text)
% ==================================================================

\subsubsection*{\normalfont\textbf{15.5.3 \quad Joint effective allowed band and working benchmark}}
\label{subsec:eps_joint_band_rebuilt}

With the five effective channels specified in
\S\ref{subsec:eps_retained_rebuilt}, we now report the resulting joint
effective allowed band, identify the working benchmark used throughout
the rebuilt \S15.5 and intended as the working reference for the
subsequent cascade updates elsewhere in the paper, and state its
relationship to the earlier conservative illustrative point used in the
previous preprint treatment of the Hubble mechanism
(\S\ref{sec:universe_age}, \S\ref{sec:mp_hubble}).

\paragraph{Headline result.}
Under the retained five-probe set, within the uniform-$\varepsilon$
effective ansatz of \S\ref{subsec:eps_def_rebuilt}, under the physical
prior $\varepsilon \geq 0$, and within the present methodology of
\S\ref{subsec:eps_retained_rebuilt}, the intersection of the $1\sigma$
intervals of the two Primary channels with the $1\sigma$ upper bounds
of the three Consistency channels defines the joint effective band
\begin{equation}
  \varepsilon \in [+0.029,\,+0.036]\ (1\sigma),
  \qquad
  \varepsilon \in [+0.021,\,+0.042]\ (2\sigma).
  \label{eq:eps_joint_band_rebuilt}
\end{equation}
The $1\sigma$ band is bounded below by the JWST channel (lower edge
$\varepsilon = 0.029$, eq.~\eqref{eq:eps_jwst_retained}) and bounded
above by the $f\sigma_8$ channel (upper edge $\varepsilon = 0.036$,
eq.~\eqref{eq:eps_fsigma8_retained}). Equivalently, the two Primary
channels alone intersect at $[+0.029,\,+0.038]$ at $1\sigma$, and the
retained Consistency channels reduce the upper edge to $+0.036$ within
the present methodology. The arithmetic midpoint of the
retained $1\sigma$ band is $\varepsilon \approx 0.032$; this number is
used only when a single illustrative value is explicitly required.

\paragraph{No-prediction caveat.}
This retained-five-probe effective band is an empirical result of the
present diagnostic analysis, not a first-principles ECT prediction. It
does not represent a measurement of a fundamental parameter of the
theory, nor a direct reconstruction of an epoch-dependent
$\varepsilon(z)$. It is an effective consistency window under the
uniform-$\varepsilon$ ansatz, the $\Lambda$CDM-background proxy, and
the retention/likelihood choices of \S\ref{subsec:eps_retained_rebuilt}.

\paragraph{Why the band is nevertheless informative.}
Three qualitative observations make the band a non-trivial consistency
check rather than a circular restatement of the retention choice.
\emph{Distinct systematics:} the five retained channels have
qualitatively distinct observational systematics, from CMB-era
distance-ratio proxies to low-$z$ RSD growth likelihoods; a narrow
common admissible region under a single uniform-$\varepsilon$ ansatz
is not guaranteed a priori. \emph{Two-sided primaries:} the two
Primary channels (Hubble$+r_s$ and JWST early-galaxy excess) provide
two-sided effective extraction intervals whose $1\sigma$ intersection
alone already defines a narrow positive-$\varepsilon$ region compatible
with the Consistency upper bounds; the band is not generated by the
$\varepsilon \geq 0$ prior alone. \emph{Narrow width:} the retained
$1\sigma$ window, $\Delta\varepsilon \approx 0.007$, is substantially
narrower than the broadest individual consistency bounds and is
therefore not generated by the least constraining channel alone.

\paragraph{Working benchmark (Benchmark B).}
The working cosmological benchmark of the present analysis is the
retained five-probe effective band itself, which we refer to as
\emph{Benchmark B}. In this sense, Benchmark B is not a single number
but the interval in Eq.~\eqref{eq:eps_joint_band_rebuilt}. When a
one-number illustration is explicitly useful, the arithmetic midpoint
of the retained $1\sigma$ band, $\varepsilon \approx 0.032$, may be
quoted as an illustration only. It should not be interpreted as a
first-principles ECT prediction or as a separately defined benchmark
constant.

\paragraph{Parameter scaling under the retained band.}
Table~\ref{tab:eps_scaling_benchmark_B} summarises the
parameterization-dependent ranges of the key quantities across cosmic
history, induced directly by the retained $1\sigma$ effective band.

\begin{table}[htbp]
\centering
\renewcommand{\arraystretch}{1.25}
\small
\begin{tabular}{@{}lccc@{}}
\toprule
\textbf{Quantity} & \textbf{$z = 0$} & \textbf{$z = 10$} & \textbf{CMB-era proxy scale} \\
\midrule
$G_{\rm eff}/G_N$ under the retained band
  & $1.000$ & $[1.15,\,1.19]$ & $[1.48,\,1.63]$ \\
$(\alpha - \beta)(z)/(\alpha - \beta)_0$ under the retained band
  & $1.000$ & $[0.918,\,0.934]$ & $[0.782,\,0.821]$ \\
\bottomrule
\end{tabular}
\caption{Parameterization-dependent ranges induced by the retained
$1\sigma$ effective band
(Eq.~\eqref{eq:eps_joint_band_rebuilt}). The entries are not direct
observables and should not be read as first-principles ECT
predictions.}
\label{tab:eps_scaling_benchmark_B}
\end{table}

\paragraph{Benchmark A (legacy retention).}
Benchmark A is retained only as a legacy conservative illustrative
point from the earlier treatment and should not be read as part of
the retained five-probe effective result. It lies outside the retained
$1\sigma$ band and is therefore not representative of the present
effective result. Its detailed closure-level numerics and robustness
checks are preserved for continuity in the legacy discussion below and
are scheduled to be migrated to the methodology appendix
(Appendix~X2 in the rebuilt structure) at a later rebuild stage.

\paragraph{Figure.}
A compact graphical summary of the retained five-probe intervals and
the resulting joint effective band will be added in the final
figure-generation pass, after regeneration under the final retained
set and the grayscale figure rules of the preprint.

\paragraph{Bridge.}
The probes not retained in the present effective
$\varepsilon$-analysis, whether for methodological or sectoral
reasons, are discussed in \S15.5.4.

% ==================================================================
% END §15.5.3 (rebuilt). §15.5.4--15.5.6 will be inserted here in
% subsequent steps. DO NOT REMOVE existing §15.5 below.
% ==================================================================
